% year: תשפ"ה
% semester: א
% moed: לדוגמה
% number: 3
% points: 16
% categories: taylor
% source: exams/מבחנים/תשפו/exam_example 2025.pdf

\begin{question}
(16 נק') העריכו בעזרת קירוב לינארי (פולינום טיילור ממעלה ראשונה) את $\sqrt[4]{18}$ ותנו חסם לשגיאה.
\end{question}

\begin{hint}
פתחו את $f(x) = \sqrt[4]{x}$ סביב $a = 16$, כי $\sqrt[4]{16} = 2$, וחסמו את שארית לגרנז' $R_1$.
\end{hint}

\begin{solution}
נגדיר $f(x) = \sqrt[4]{x} = x^{1/4}$ ונפתח סביב $a = 16$, הנקודה הקרובה ל-$18$ שבה השורש הרביעי ידוע. עבור $x > 0$:
\[ f'(x) = \frac14 x^{-3/4}, \qquad f''(x) = -\frac{3}{16} x^{-7/4}, \]
ולכן $f(16) = 2$ ו-$f'(16) = \frac{1}{4\cdot 16^{3/4}} = \frac{1}{4\cdot 8} = \frac{1}{32}$.

\textbf{קירוב לינארי:}
\[ P_1(x) = f(16) + f'(16)(x - 16) = 2 + \frac{x - 16}{32}, \qquad
   \sqrt[4]{18} \approx P_1(18) = 2 + \frac{2}{32} = 2 + \frac{1}{16} = \frac{33}{16} = 2.0625. \]

\textbf{חסם לשגיאה.} לפי משפט טיילור עם שארית לגרנז', קיימת $c \in (16, 18)$ כך ש-
\[ R_1(18) = \frac{f''(c)}{2!}(18 - 16)^2 = \frac{-\frac{3}{16}c^{-7/4}}{2}\cdot 4 = -\frac{3}{8\,c^{7/4}}. \]
מכיוון ש-$c > 16$, מתקיים $c^{7/4} > 16^{7/4} = 2^7 = 128$, ולכן
\[ |R_1(18)| = \frac{3}{8c^{7/4}} < \frac{3}{8\cdot128} = \frac{3}{1024} \approx 0.0029. \]

\textbf{תשובה:} $\sqrt[4]{18} \approx \frac{33}{16} = 2.0625$, עם שגיאה קטנה מ-$\frac{3}{1024}$. מכיוון ש-$R_1 < 0$, הקירוב גדול מהערך האמיתי: $\frac{33}{16} - \frac{3}{1024} < \sqrt[4]{18} < \frac{33}{16}$.

(בדיקה: $\sqrt[4]{18} \approx 2.05977$, והשגיאה בפועל כ-$0.00273 < \frac{3}{1024}$.)
\end{solution}
